\chapter{Kleinian groups}

\section{Limit sets and discontinuity}

\begin{definition}[Limit set and elementary group]
\label{thurston-ch08-limit-set}
\group{Kleinian groups}\source{thurston:page-0186}\notready
For a discrete group $\Gamma$ of orientation-preserving isometries of $H^n$, the limit set $L_\Gamma\subset S^{n-1}_\infty$ is the accumulation set of any orbit $\Gamma x$.  It is independent of $x$.  The group is elementary when $\abs{L_\Gamma}\leq2$.
\end{definition}

\begin{proposition}[Minimality of the non-elementary limit set]
\label{thurston-ch08-limit-set-minimality}
\dcref{Proposition 8.1.2}\group{Kleinian groups}\source{thurston:page-0186,thurston:page-0187}\notready
If $\Gamma$ is non-elementary, every nonempty closed $\Gamma$-invariant subset of $S^{n-1}_\infty$ contains $L_\Gamma$.
\end{proposition}
\begin{proof}\sketch Take the hyperbolic convex hull of the invariant closed set.  It is invariant, and the limit points of an orbit in its interior lie in the set; hence minimality follows from convexity.\end{proof}

\begin{proposition}[Elementary-group characterization]
\label{thurston-ch08-elementary-abelian}
\dcref{Proposition 8.1.1}\group{Kleinian groups}\source{thurston:page-0186}\notready
\Gamma$ is elementary if and only if it has an abelian subgroup of finite index.
\end{proposition}

\begin{corollary}[Limit set of a normal subgroup]
\label{thurston-ch08-normal-subgroup-limit-set}
\dcref{Corollary 8.1.3}\group{Kleinian groups}\source{thurston:page-0187}\notready
If $\Gamma$ is non-elementary and $1\neq\Gamma'\triangleleft\Gamma$, then $L_{\Gamma'}=L_\Gamma$.
\end{corollary}
\begin{proof}\sketch Normality makes $L_{\Gamma'}$ invariant under $\Gamma$.  The subgroup is infinite, so minimality gives $L_\Gamma\subseteq L_{\Gamma'}$; the reverse inclusion follows from orbit containment.\end{proof}

\begin{definition}[Domain of discontinuity and Kleinian group]
\label{thurston-ch08-domain-discontinuity}
\group{Kleinian groups}\source{thurston:page-0189}\notready
Set $D_\Gamma=S^{n-1}_\infty\setminus L_\Gamma$.  A discrete subgroup of $\PSL(2,\C)$ with $D_\Gamma\neq\varnothing$ is called a Kleinian group.
\end{definition}

\begin{definition}[Properly discontinuous action]
\label{thurston-ch08-proper-discontinuity}
\dcref{Definition 8.2.1}\group{group actions}\source{thurston:page-0189}\notready
An action of a group on a locally compact space $X$ is properly discontinuous when, for every compact $K\subset X$, only finitely many $\gamma$ satisfy $\gamma K\cap K\neq\varnothing$.
\end{definition}

\begin{proposition}[Quotients of proper actions]
\label{thurston-ch08-proper-quotient}
\dcref{Proposition 8.2.2}\group{group actions}\source{thurston:page-0189}\notready
For a properly discontinuous action on a locally compact Hausdorff space, $X/\Gamma$ is Hausdorff; if the action is free, $X\to X/\Gamma$ is a covering projection.
\end{proposition}
\begin{proof}\sketch Separate distinct orbits using compact neighborhoods and finitely many intersecting translates.  Freeness permits a neighborhood disjoint from all its translates, which projects evenly.\end{proof}

\begin{proposition}[Discontinuity on the domain]
\label{thurston-ch08-domain-proper}
\dcref{Proposition 8.2.3}\group{Kleinian groups}\source{thurston:page-0189,thurston:page-0190}\notready
The action of a discrete isometry group $\Gamma$ on $D_\Gamma$ and on $H^n\cup D_\Gamma$ is properly discontinuous.
\end{proposition}
\begin{proof}\sketch Retract $H^n\cup D_\Gamma$ to the convex hull $H(L_\Gamma)$ by nearest-point or horospherical projection.  The retraction is proper, and discreteness gives proper discontinuity on the convex hull minus its ideal boundary.\end{proof}

\begin{definition}[Wandering action and discrete orbits]
\label{thurston-ch08-wandering-discrete-orbits}
\group{group actions}\source{thurston:page-0190,thurston:page-0191}\notready
An action is wandering if each point has a neighborhood meeting only finitely many of its translates.  It has discrete orbits if every orbit has empty limit set.
\end{definition}

\begin{proposition}[Wandering free quotients]
\label{thurston-ch08-wandering-covering}
\dcref{Proposition 8.2.6}\group{group actions}\source{thurston:page-0190}\notready
The quotient map of a free wandering action on a Hausdorff space is a covering projection.  Wandering alone need not make the quotient Hausdorff.
\end{proposition}

\section{Convex hyperbolic manifolds}

\begin{definition}[Convex hyperbolic manifold]
\label{thurston-ch08-convex-manifold}
\dcref{Definition 8.3.1}\group{convex hyperbolic manifolds}\source{thurston:page-0192}\notready
A complete hyperbolic manifold with boundary is convex if every path is homotopic rel endpoints to a geodesic arc.
\end{definition}

\begin{proposition}[Developing-map criterion]
\label{thurston-ch08-convex-developing}
\dcref{Proposition 8.3.2}\group{convex hyperbolic manifolds}\source{thurston:page-0192}\notready
A complete hyperbolic manifold is convex exactly when its developing map identifies its universal cover homeomorphically with a convex subset of $H^n$.
\end{proposition}
\begin{proof}\sketch Convexity of the developing image straightens lifted paths.  Conversely, path straightening proves injectivity of the developing map and convexity of its image.\end{proof}

\begin{proposition}[Local convexity and deformation stability]
\label{thurston-ch08-convex-deformation}
\dcref{Proposition 8.3.3}\group{convex hyperbolic manifolds}\source{thurston:page-0192,thurston:page-0193,thurston:page-0194}\notready
Local convexity implies convexity.  Every sufficiently small deformation of a compact convex hyperbolic manifold can be enlarged slightly to a convex hyperbolic manifold homeomorphic to it.  The same holds for uniformly convex noncompact manifolds under uniformly small deformations.
\end{proposition}
\begin{proof}\sketch Subdivide paths into uniformly short pieces and straighten adjacent pieces; strict convex enlargement supplies a uniform interior buffer.  For a deformation, take convex hulls in uniformly sized balls; their union remains locally convex and hence convex.\end{proof}

\begin{proposition}[Convex-core boundary correspondence]
\label{thurston-ch08-convex-quasiconformal}
\dcref{Proposition 8.3.4}\group{convex hyperbolic manifolds}\source{thurston:page-0194}\notready
If $M_1,M_2$ are strictly convex compact hyperbolic $n$-manifolds and a homotopy equivalence $M_1\to M_2$ is a boundary diffeomorphism, then the induced groups on the disk are conjugate by a quasiconformal homeomorphism of $B^n$ that is a pseudo-isometry on $H^n$.
\end{proposition}
\begin{proof}\sketch Lift the equivalence and extend it along corresponding normal rays.  Comparison of distances between nearby normal rays gives the pseudo-isometry and its quasiconformal boundary extension.\end{proof}

\begin{definition}[Kleinian quotient spaces]
\label{thurston-ch08-kleinian-quotients}
\group{Kleinian groups}\source{thurston:page-0194,thurston:page-0195}\notready
Write
\[M_\Gamma=H(L_\Gamma)/\Gamma,\quad N_\Gamma=H^n/\Gamma,\quad O_\Gamma=(H^n\cup D_\Gamma)/\Gamma,\quad P_\Gamma=(H^n\cup D_\Gamma\cup W_\Gamma)/\Gamma,\]
where $W_\Gamma$ is the projective set dual to planes whose ideal boundary lies in $D_\Gamma$.  There are inclusions $M_\Gamma\subset N_\Gamma\subset O_\Gamma\subset P_\Gamma$; the first three have the same homotopy type, $N_\Gamma=\Int(O_\Gamma)$, and $M_\Gamma,O_\Gamma$ are homeomorphic outside degenerate cases.
\end{definition}

\section{Geometric finiteness}

\begin{definition}[Geometrically finite group]
\label{thurston-ch08-geometrically-finite}
\dcref{Definition 8.4.1}\group{geometric finiteness}\source{thurston:page-0195}\notready
\Gamma$ is geometrically finite if a fixed-radius neighborhood of $M_\Gamma$ has finite volume.
\end{definition}

\begin{theorem}[Ahlfors dichotomy]
\label{thurston-ch08-ahlfors-dichotomy}
\dcref{Theorem 8.4.2}\group{geometric finiteness}\source{thurston:page-0195,thurston:page-0196}\notready
For a geometrically finite $\Gamma$, $L_\Gamma$ has spherical measure zero or full measure.  In the full-measure case the action on $S^2_\infty$ is ergodic.
\end{theorem}
\begin{proof}\sketch Extend an invariant bounded boundary function harmonically by visual averaging.  The gradient flow preserves volume; outside the convex core the average is strictly below one half, while a finite-volume superlevel set cannot support a complete volume-preserving flow.  Density points then force the boundary function to be almost everywhere constant.\end{proof}

\begin{proposition}[Equivalent finiteness criteria]
\label{thurston-ch08-geometric-finiteness-equivalences}
\dcref{Proposition 8.4.3}\group{geometric finiteness}\source{thurston:page-0197,thurston:page-0198,thurston:page-0199}\notready
For a discrete torsion-free subgroup of $\PSL(2,\C)$, the following are equivalent: $\Gamma$ is geometrically finite; the thick part $M_{[\epsilon,\infty)}$ is compact; and $\Gamma$ has a finite-sided fundamental polyhedron.
\end{proposition}
\begin{proof}\sketch Finite volume gives compactness by disjoint embedded balls.  Compact thick part plus standard cusp strips yields finitely many-sided cusp and core polyhedra.  Conversely, tangencies of translates of a finite polyhedron can occur only at finitely many parabolic cusps, where finite-width strips give finite-volume neighborhoods.\end{proof}

\section{Convex-hull boundary and quasi-Fuchsian groups}

\begin{proposition}[Convex-hull boundary surface]
\label{thurston-ch08-convex-boundary-surface}
\dcref{Proposition 8.5.1}\group{convex geometry}\source{thurston:page-0200}\notready
For torsion-free Kleinian $\Gamma$, $\partial M_\Gamma$ has a complete hyperbolic surface structure; its bending locus is a closed geodesic lamination.
\end{proposition}

\begin{definition}[Lamination]
\label{thurston-ch08-lamination}
\dcref{Definition 8.5.2}\group{laminations}\source{thurston:page-0200,thurston:page-0201}\notready
A lamination is a closed subset locally modeled on $\R^{n-k}\times B$ with coordinate changes preserving the transverse $B$-coordinate.  Its leaves are the connected leafwise pieces; a geodesic lamination has geodesic leaves.
\end{definition}

\begin{definition}[Transverse measure]
\label{thurston-ch08-transverse-measure}
\dcref{Definition 8.6}\group{measured laminations}\source{thurston:page-0204}\notready
For a lamination with local transverse spaces $B_i$, a transverse measure is a measure on each $B_i$, finite on compact transversals, invariant under the transverse coordinate changes.  Equivalently it is a measure on every transversal, supported on the lamination and invariant under projection along leaves.
\end{definition}

\begin{proposition}[Convex-hull bending measure]
\label{thurston-ch08-bending-measure}
\group{measured laminations}\source{thurston:page-0204,thurston:page-0205,thurston:page-0206}\notready
The bending locus of $\partial M_\Gamma$ carries a canonical transverse measure $\beta$: on a transverse arc, $\beta$ is the total turning angle of the outward normal.  It is obtained as the monotone limit of outer polyhedral approximations.  The hyperbolic surface structure of $\partial M_\Gamma$, its bending lamination, and $\beta$ determine the local isometric embedding near the boundary.  Every isolated leaf in the support is a closed geodesic, and for three pairwise intersecting planes the dihedral angles satisfy $\alpha+\beta\leq\gamma$.
\end{proposition}
\begin{proof}\sketch Approximate the convex boundary by finite unions of planes.  The triangle-area inequality gives monotonicity of angle sums, so their limit defines a transverse measure; short-arc limits identify it with normal turning and determine the local embedding.
\end{proof}

\begin{definition}[Quasi-Fuchsian group]
\label{thurston-ch08-quasi-fuchsian}
\dcref{Definition 8.7.1}\group{quasi-Fuchsian groups}\source{thurston:page-0206,thurston:page-0207}\notready
A Kleinian group is quasi-Fuchsian if its limit set is a Jordan curve in $S^2_\infty$ (equivalently, its domain of discontinuity has two invariant components).
\end{definition}

\begin{proposition}[Marden quasi-Fuchsian criterion]
\label{thurston-ch08-marden-criterion}
\dcref{Proposition 8.7.2}\group{quasi-Fuchsian groups}\source{thurston:page-0207}\notready
For a torsion-free Kleinian group, quasi-Fuchsianity is equivalent to $N_\Gamma$ being homeomorphic to $S\times\R$ with compact convex core and to each component of $\partial M_\Gamma$ being a finite-area surface whose inclusion is a homotopy equivalence.
\end{proposition}
\begin{proof}\sketch A quasiconformal conjugacy transfers the Fuchsian product picture.  Conversely, the two boundary components and the convex deformation theorem give a quasiconformal sphere conjugacy, whose invariant separating curve is the limit set.\end{proof}

\begin{definition}[Complete geodesic lamination]
\label{thurston-ch08-complete-lamination}
\dcref{Definition 8.7.5}\group{geodesic laminations}\source{thurston:page-0211}\notready
A geodesic lamination on a hyperbolic surface is complete when every complementary component is an ideal polygon (with the finite-area cusp version understood).
\end{definition}

\begin{proposition}[Geodesic completion]
\label{thurston-ch08-geodesic-lamination-completion}
\dcref{Proposition 8.7.6}\group{geodesic laminations}\source{thurston:page-0211}\notready
Every geodesic lamination is contained in a complete geodesic lamination.
\end{proposition}

\section{Uncrumpled surfaces and compactness}

\begin{definition}[Uncrumpled surface and wrinkling locus]
\label{thurston-ch08-uncrumpled-surface}
\dcref{Definition 8.8.1}\group{uncrumpled surfaces}\source{thurston:page-0215}\notready
An uncrumpled surface in a hyperbolic three-manifold is a complete finite-area hyperbolic surface with an isometric map $f:S\to N$ such that every point lies in the interior of a mapped straight segment and cusps map to cusps.  Its wrinkling locus is the lamination where the map fails to be locally totally geodesic.
\end{definition}

\begin{proposition}[Wrinkling lamination]
\label{thurston-ch08-wrinkling-lamination}
\dcref{Proposition 8.8.2}\group{uncrumpled surfaces}\source{thurston:page-0215}\notready
The wrinkling locus of an uncrumpled surface is a geodesic lamination, and the map is totally geodesic off it.
\end{proposition}

\begin{proposition}[Mumford compactness and surface compactness]
\label{thurston-ch08-uncrumpled-compactness}
\dcref{Proposition 8.8.3,Proposition 8.8.4,Proposition 8.8.5}\group{uncrumpled surfaces}\source{thurston:page-0216,thurston:page-0217,thurston:page-0218}\notready
For each topological surface $S$ and $\epsilon>0$, the $\epsilon$-thick part of moduli space is compact.  There is a uniform $\epsilon$ below which every short geodesic is separating in the required cusp-controlled sense.  For a compact $K\subset N$, the uncrumpled surfaces of fixed type intersecting $K$ and with a uniform lower bound on null-homotopic loop lengths form a compact set.
\end{proposition}
\begin{proof}\sketch Choose bounded-area fundamental domains, use thick-part compactness to prevent degeneration, and apply geometric convergence plus the straightening condition to obtain a limiting uncrumpled surface.\end{proof}

\begin{definition}[Train track]
\label{thurston-ch08-train-track}
\dcref{Definitions 8.9.1}\group{train tracks}\source{thurston:page-0220}\notready
A train track on a differentiable surface is a smoothly embedded branched one-complex whose branches meet at switches with common tangent directions; its complementary regions satisfy the standard negative-Euler-characteristic conditions.  A lamination is carried when it lies in the associated foliated neighborhood and follows its branch foliation.
\end{definition}

\begin{proposition}[Train-track approximation]
\label{thurston-ch08-train-track-approximation}
\dcref{Proposition 8.9.2,Corollary 8.9.3}\group{train tracks}\source{thurston:page-0221}\notready
For every hyperbolic surface $S$ and $\delta>0$ there is $\epsilon>0$ such that every geodesic lamination has a train-track approximation with branch curvature $<\delta$.  Laminations carried by a sufficiently close approximation have leaves uniformly close to those of the original lamination.
\end{proposition}
\begin{proof}\sketch Use short leaves of the complementary foliation to build switches tangent to the geodesic lamination, grouping short branches when necessary.  Small curvature gives uniform closeness to the unique geodesic with the same endpoints.\end{proof}

\begin{proposition}[Geodesic representative of a lamination]
\label{thurston-ch08-geodesic-representative}
\dcref{Proposition 8.9.4}\group{geodesic laminations}\source{thurston:page-0221}\notready
A lamination on a hyperbolic surface is isotopic to a geodesic lamination if and only if it has no Reeb component and its leaves have the prescribed non-intersection behavior; the geodesic representative is unique in its isotopy class.
\end{proposition}

\section{Realization}

\begin{definition}[Realizable lamination]
\label{thurston-ch08-realizable-lamination}
\dcref{Definition 8.10.1}\group{lamination realization}\source{thurston:page-0223}\notready
For a map $f:S\to N$ from a hyperbolic surface to a complete hyperbolic three-manifold, a geodesic lamination $\gamma$ is realizable in the homotopy class of $f$ if some homotopic uncrumpled surface has wrinkling locus containing $\gamma$.  Write $\mathcal R_f$ for the set of realizable laminations.
\end{definition}

\begin{proposition}[Realization criterion and openness]
\label{thurston-ch08-realization-openness}
\dcref{Proposition 8.10.2}\group{lamination realization}\source{thurston:page-0223,thurston:page-0224}\notready
If $\pi_1f$ is injective and preserves the relevant non-peripheral classes, every finite enlargement of a realizable lamination is realizable; $\mathcal R_f$ is open in the train-track topology.
\end{proposition}
\begin{proof}\sketch Straighten finitely many leaves in the universal cover; distinct endpoints at infinity give disjoint geodesics.  Train-track approximation and compactness of uncrumpled surfaces extend the construction to neighborhoods and arbitrary enlargements.\end{proof}

\begin{proposition}[Topology of lamination spaces]
\label{thurston-ch08-lamination-space-topology}
\dcref{Proposition 8.10.3,Proposition 8.10.4,Proposition 8.10.5,Proposition 8.10.6}\group{lamination realization}\source{thurston:page-0224,thurston:page-0225}\notready
The map from measured laminations to geodesic laminations is continuous.  The wrinkling-locus map from uncrumpled surfaces is continuous.  The space $\mathcal{GL}(S)$ of geodesic laminations is compact, the projective measured-lamination space is a compact manifold, and every geodesic lamination admits a transverse measure.
\end{proposition}

\begin{proposition}[Finite-leaf density and end behavior]
\label{thurston-ch08-finite-leaf-density}
\dcref{Proposition 8.10.7}\group{lamination realization}\source{thurston:page-0227,thurston:page-0228,thurston:page-0229}\notready
For a hyperbolic surface $S$:
\begin{enumerate}
\item the elements of $\mathcal{ML}(S)$ whose supports are finite unions of compact or proper geodesics form a dense subset;
\item finite-leaf points form a dense subset of $\mathcal{GL}(S)$;
\item each end of a noncompact leaf in such a lamination either approaches a cusp or accumulates by winding around a closed geodesic.
\end{enumerate}
\end{proposition}
\begin{proof}\sketch On a train track, positive transverse measures satisfy linear equations with integer coefficients, so rational solutions approximate them; clearing denominators produces finitely many compact or proper leaves.  For an arbitrary lamination, take a close train-track approximation and retain the branches lying on cyclic or proper train routes.  Cover those branches by finitely many such routes, then add finitely many leaves along shortest routes from each remaining branch, merging their ends by spiraling around compact leaves or following proper leaves.  Finally, an end of a noncompact leaf in a finite-leaf lamination is either proper or accumulates on a compact leaf, in which case the local order of nearby leaves forces it to spiral around that leaf.\end{proof}

\begin{theorem}[Open dense realization theorem]
\label{thurston-ch08-realization-dichotomy}
\uses{thurston-ch08-finite-leaf-density}
\dcref{Theorem 8.10.8}\group{lamination realization}\source{thurston:page-0228,thurston:page-0229,thurston:page-0230}\notready
Let $f:S\to N$ have injective fundamental-group map and preserve non-parabolicity.  Then $\mathcal R_f$ is an open dense subset of $\mathcal{GL}(S)$.
\end{theorem}
\begin{proof}
\sketch The hypotheses make every finite-leaf lamination nondegenerate and realizable, so finite-leaf density makes $\mathcal R_f$ dense.  For openness, first straighten a single added leaf whose ends are asymptotic to leaves already realized.  Finite additions follow inductively, and Proposition 8.10.7 applied to the double of $S-\gamma$ approximates an arbitrary enlargement of $\gamma$ by finite additions.  Compactness and continuity of the relevant uncrumpled surfaces pass realizability to the limit.  A finite collection of close train-track neighborhoods for those surfaces then gives a neighborhood of $\gamma$ contained in $\mathcal R_f$.
\end{proof}

\begin{corollary}[Geometrically finite realization alternative]
\label{thurston-ch08-geometrically-finite-realization}
\uses{thurston-ch08-realization-dichotomy}
\dcref{Corollary 8.10.9}\group{lamination realization}\source{thurston:page-0230,thurston:page-0231}\notready
For geometrically finite $\Gamma$ and an admissible map $f:S\to N_\Gamma$, either $\mathcal R_f=\mathcal{GL}(S)$ or a finite-index subgroup of $\Gamma$ is the group of a finite-volume manifold fibering over $S^1$.
\end{corollary}
\begin{proof}
\sketch Geometric finiteness makes the space of uncrumpled surfaces homotopic to $f$ compact.  If the group of mapping classes arising from return homotopies of these surfaces is finite, compactness of their wrinkling loci together with the open dense realization theorem gives $\mathcal R_f=\mathcal{GL}(S)$.  Otherwise, after passage to finite index, one obtains an infinite-order mapping class $\phi$ for which $f$ and $f\circ\phi$ are conjugate by some $\beta\in\Gamma$.  The group generated by $f(\pi_1S)$ and $\beta$ is then the fundamental group of the mapping torus of $\phi$, giving the finite-volume fibering alternative.
\end{proof}

\begin{conjecture}[Quasi-Fuchsian realization conjecture]
\label{thurston-ch08-realization-conjecture}
\dcref{Conjecture 8.10.10}\group{lamination realization}\source{thurston:page-0230}\notready
For any cusp-preserving map $f:S\to N$, the image subgroup is quasi-Fuchsian if and only if every geodesic lamination of $S$ is realizable in the homotopy class of $f$.
\end{conjecture}

\section{Cusps and geometric tameness}

\begin{proposition}[Parabolic curves at extra cusps]
\label{thurston-ch08-extra-cusps}
\dcref{Proposition 8.11.1}\group{cusps}\source{thurston:page-0231,thurston:page-0232}\notready
If $f:S\to N$ is a relative homotopy equivalence and $N$ has parabolics not represented by cusps of $S$, then each end has a finite geodesic lamination $\gamma_+$ or $\gamma_-$ of disjoint simple closed curves, in bijection with its extra cusps.  An element of $\pi_1S$ maps to a parabolic exactly when it is peripheral on $S$ or freely homotopic to a leaf of one of these laminations.
\end{proposition}
\begin{proof}\sketch Represent extra cusps by half-open cylinders ending on an incompressible separating surface.  Homotope their boundary curves to geodesics; two intersections would give lifts whose geodesic boundary lines meet twice, impossible.\end{proof}

\begin{definition}[Geometrically tame end]
\label{thurston-ch08-geometrically-tame-end}
\dcref{Definition 8.11.2}\group{geometric tameness}\source{thurston:page-0233,thurston:page-0234}\notready
Let $P$ be cusp horoball neighborhoods and $M$ the convex hull.  An end $E$ of $N-P$ with incompressible boundary surface $S\subset M$ is geometrically tame if either $E\cap M$ is compact or the family of uncrumpled surfaces homotopic to $S$ and eventually contained in $E$ is noncompact.  A manifold is geometrically tame when it has a compact core whose complementary ends are geometrically tame.
\end{definition}

\section{Harmonic functions and ergodicity}

\begin{proposition}[Flow-time estimate]
\label{thurston-ch08-flow-time-estimate}
\dcref{Equation 8.12.1}\group{harmonic functions}\source{thurston:page-0234}\notready
For a positive function $h$ and the flow generated by $-\operatorname{grad}h$, if $A$ is a finite-length portion of a flow line through $x$, then
\[
T(A)=\int_A\frac{1}{\norm{\operatorname{grad}h}}\,ds\geq\frac{\length(A)^2}{\int_A\norm{\operatorname{grad}h}\,ds}\geq\frac{\length(A)^2}{h(x)}.
\]
The flow exists for all positive time.
\end{proposition}

\begin{proposition}[Neighborhood-volume estimate]
\label{thurston-ch08-uncrumpled-neighborhood-volume}
\dcref{Proposition 8.12.1}\group{harmonic functions}\source{thurston:page-0234,thurston:page-0235}\notready
If an uncrumpled surface $S\xrightarrow f N$ has no null-homotopic loop of length at most $a$, then $\vol(\mathcal N_1(f(S)))\leq-C\chi(S)$ for a constant $C$ depending only on $a$.
\end{proposition}
\begin{proof}\sketch Integrate over $S$ the multiplicity of short path classes into an auxiliary function on $N$.  Its total integral is bounded by area$(S)$ times the volume of a fixed ball, while on the one-neighborhood it is bounded below by the area of a radius-$a/2$ disk.\end{proof}

\begin{theorem}[Boundary minimum principle for tame manifolds]
\label{thurston-ch08-tame-harmonic-minimum}
\dcref{Theorem 8.12.3}\group{harmonic functions}\source{thurston:page-0235,thurston:page-0236,thurston:page-0237}\notready
If $N$ is geometrically tame and $M$ is its convex hull, every nonconstant positive harmonic function on $M$ satisfies $\inf_Mh=\inf_{\partial M}h$.  The same conclusion holds for positive superharmonic functions.
\end{theorem}
\begin{proof}\sketch Exhaust the tame ends by separated uncrumpled surfaces of uniformly bounded one-neighborhood volume.  Their modified cusp cycles bound chains whose supports exhaust $M$.  The flow-time estimate forces any positive-measure family of gradient flow lines to spend quadratically growing time in sets of only linearly growing volume unless it exits through $\partial M$; hence no interior minimum can be lower than the boundary infimum.\end{proof}

\begin{corollary}[Limit-set ergodicity for tame groups]
\label{thurston-ch08-tame-limit-ergodicity}
\dcref{Corollary 8.12.4}\group{harmonic functions}\source{thurston:page-0236}\notready
If $N$ is geometrically tame and $\Gamma=\pi_1N$, then $L_\Gamma$ has measure zero or one; in the latter case $\Gamma$ acts ergodically on $S^2_\infty$.
\end{corollary}
\begin{proof}\sketch Harmonically extend the characteristic function of an invariant measurable subset of $L_\Gamma$.  The boundary minimum principle makes this extension constant when $L_\Gamma=S^2$ and forces the subset to have measure zero otherwise.\end{proof}
